When a program changes, a differential check that finds no difference is easily read as reassurance, although it only reports what the chosen inputs exercised. We describe V7-B, a prototype that compares two versions of a Python function by executing both, and that emits exactly one of three verdicts: \DIV{} backed by a replayable witness, \INS{} with a coded reason, or \NDO{} with a recorded sufficiency level (S1 or S2) that states how much of the changed code was exercised. No equivalence verdict exists. Evidence is stored as content-addressed, hash-chained records so that a verdict can be traced to the runs behind it. We evaluate the prototype only exploratorily on 321 development pairs, all of which are exposed or non-pristine (QuixBugs, BugsInPy function units, synthetic mutants). On these data V7-B detects 0.708 of the 106 pairs with a divergent label, against 0.774 for plain same-input differential execution (B1); adding the same repetition counts to B1 (B1+rep) changes the verdict on only 2 of 321 pairs, and the sufficiency level does not visibly separate generator labels. V7-B abstains on 0.360 of 258 primary pairs, mostly because of unimplemented engine features rather than a meaningful sufficiency judgement. Power on the observed design is too low for a confirmatory claim, the real-pair feasibility gate K1 is not met, and no published tool was compared. We therefore present V7-B as a design and feasibility report with negative and null findings, not as a validation.
